% =============================================================================
% Glossar der Komponentennamen
% =============================================================================

\section*{Glossar}
\addcontentsline{toc}{section}{Glossar}
\label{sec:glossar}

Dieses Glossar erklärt die Prozess- und Komponentennamen des ROSI-Systems, die in
Kapitel~\ref{sec:architektur} eingeführt und im weiteren Text ohne erneute Erläuterung
verwendet werden.

\begin{table}[htbp]
    \centering
    \renewcommand{\arraystretch}{1.3}
    \begin{tabularx}{\textwidth}{lX}
        \hline
        \textbf{Begriff} & \textbf{Erklärung} \\
        \hline
        FrameGrabber & Service-Prozess, der die Kamerahardware anbindet, rohe Pixeldaten
        in Shared-Memory-Blöcke schreibt und Referenzen darauf an den PipelineManager
        weiterreicht. \\
        HealthMonitor & Hintergrund-Thread im MPWorkerPool, der Heartbeats der
        Worker-Prozesse überwacht und einen nicht mehr antwortenden Worker neu startet. \\
        InspectionDataHandler & Service-Prozess, der fortlaufend leere Daten-Container
        erzeugt und über eine begrenzte Queue bereitstellt; wirkt bei voller Queue als
        Backpressure-Ventil gegen unbegrenzte Speicherallokation. \\
        LoggerProcess & Service-Prozess, der Lognachrichten aller übrigen Prozesse und
        Worker über eine zentrale Queue entgegennimmt und asynchron persistiert. \\
        main.py & Startprozess von ROSI; einziger Supervisor im System, überwacht zur
        Laufzeit ausschließlich den PipelineManager. \\
        MPWorkerPool & Pool aus zwölf Worker-Prozessen, der rechenintensive
        Pipeline-Schritte über ein Map-Reduce-Muster parallel verarbeitet. \\
        PipelineManager & Orchestriert die sechsstufige Verarbeitungspipeline jeder
        Inspektion; einziger vom Hauptprozess laufzeitüberwachter Service-Prozess. \\
        ResultCollector & Hintergrund-Thread im MPWorkerPool, der Teilergebnisse der
        Worker nicht blockierend einsammelt und die Pipeline nach Abschluss aller
        Teiljobs benachrichtigt. \\
        \hline
    \end{tabularx}
\end{table}
